% 07-bai-tap.tex - End-of-chapter exercises (three tiers)
%
% Scope: tier 1 (understand), tier 2 (apply/code), tier 3 (extend/open).

\exercises

\index{bài tập!kiểm tra hiểu|(}%

\par\textbf{Hướng dẫn chung:} Code cho toàn bộ chương này nằm trong companion
repo \texttt{rnn-to-transformer-lab} tại tag \texttt{ch01}. Clone repo, checkout
tag, cài môi trường conda, rồi mở file
\mintinline{text}{rnn_to_transformer_lab/__init__.py}. Đọc code cùng với mục~4
và~5 của chương. Khi một bài yêu cầu bạn sửa code, hãy sửa trực tiếp file đó
và chạy \mintinline{text}{python verify.py} để kiểm tra.

\index{bài tập!kiểm tra hiểu|)}%

\tierunderstand

\begin{bt}
  \item Tại bước $t$, $\vect{h}_t$ nhận gradient từ hai nguồn. Kể tên hai
    nguồn đó và viết công thức Jacobian cho từng nguồn. Nguồn nào biến mất
    tại bước cuối cùng $t = T$?

  \item Mở companion repo, chạy \mintinline{text}{python verify.py}, và đọc
    kết quả lan truyền xuôi. Không cần tính toán: nhìn vào $\vect{h}_3$ và
    giải thích vì sao các phần tử của nó đều gần bằng 1. Gợi ý: nhìn vào
    $\vect{a}_3$.

  \item Trong BPTT, gradient qua $\tanh$ là phép nhân Hadamard với
    $(\vect{1} - \vect{h}_t^2)$. Một phần tử của $\vect{h}_t$ có giá trị
    $0.95$. Hỏi: tín hiệu gradient qua phần tử đó bị suy giảm bao nhiêu lần
    so với một phần tử có giá trị $0.1$?

  \item Nhìn vào $\partial L / \partial \mat{W}_{hh}$ trong mục~5. Phần tử
    nào lớn nhất, và nó nối hidden unit nào với hidden unit nào? Lần ngược
    về lan truyền xuôi để giải thích vì sao.

  \item Vẽ đồ thị tính toán tháo cuộn của RNN ba bước. Trên mỗi cạnh, ghi
    kích thước của Jacobian tương ứng. Mũi tên nào là nơi $\mat{W}_{hh}$ xuất
    hiện nhiều lần nhất?
\end{bt}

\tierapply

\index{bài tập!để tay vào|(}%

\begin{bt}
  \item Trong companion repo, đổi \tn{hàm kích hoạt}{activation function} từ
    $\tanh$ sang
    \mintinline{python}{relu} (dùng \mintinline{python}{np.maximum(0, a)}).
    Chạy lại \mintinline{python}{python verify.py}. Gradient check còn khớp
    không? Nếu không, vì sao? Gợi ý: ReLU không khả vi tại 0, nhưng
    \tn{sai phân hữu hạn}{finite difference} thì vẫn chạy.

  \item Thêm một \tn{bước thời gian}{timestep} thứ tư với
    $\vect{x}_4 = [0.2, -0.5]$ và
    $\widehat{\vect{y}}_4 = [0.5, 0.0]$. Cập nhật \mintinline{python}{X} và
    \mintinline{python}{TARGET} trong code. Chạy gradient check. Gradient
    của $\mat{W}_{hh}$ thay đổi thế nào so với bản ba bước? Giải thích.

  \item Cố ý gây một lỗi BPTT: trong hàm \mintinline{python}{backward}, bỏ
    \mintinline{python}{dh_next} khỏi dòng tính \mintinline{python}{dh_t}
    (tức là chỉ giữ đường gradient từ output). Gradient check bắt được lỗi
    này không? Sai số tương đối tăng lên bao nhiêu?

  \item Cài đặt BPTT cho bài toán phân loại: thay vì hồi quy, đặt một softmax
    lên $\vect{y}_T$ (chỉ bước cuối) và dùng cross-entropy. Đạo hàm qua
    softmax + cross-entropy là $\vect{y}_T - \widehat{\vect{y}}_T$, giống hệt
    trường hợp hồi quy; chỉ khác là nó chỉ áp dụng ở bước cuối, và
    $\widehat{\vect{y}}_T$ là one-hot. Cài đặt và kiểm tra gradient.
\end{bt}

\index{bài tập!để tay vào|)}%

\tierextend

\index{bài tập!đẩy xa hơn|(}%

\begin{bt}
  \item Đo thời gian chạy của \mintinline{python}{backward()} và
    \mintinline{python}{finite_difference()} cho $T = 3, 10, 50, 100$.
    Vẽ đồ thị. Với $T = 100$, sai phân hữu hạn chậm hơn BPTT bao nhiêu lần?
    Từ đó giải thích vì sao mọi framework deep learning đều dùng
    \tn{lan truyền ngược}{backpropagation} thay vì sai phân hữu hạn, và vì sao
    sai phân hữu hạn vẫn là công cụ
    kiểm tra không thể thiếu.

  \item Đọc bài báo ``Automatic Differentiation in Machine Learning: a
    Survey'' của Baydin và cộng sự (2018, arXiv:1502.05767). So sánh cách
    bài báo đó trình bày forward-mode AD và reverse-mode AD (BPTT là
    reverse-mode) với dẫn xuất Jacobian trong chương này. Viết một đoạn ngắn
    giải thích vì sao reverse-mode có lợi thế khi số tham số lớn hơn số đầu
    ra.

  \item Dùng thư viện \mintinline{python}{torch} (PyTorch) cài đặt lại chính
    xác RNN ba bước này, nhưng để cơ chế autograd của PyTorch tự tính
    gradient. So sánh gradient của PyTorch với gradient BPTT thủ công của bạn
    (gợi ý: dùng \mintinline{python}{torch.from_numpy} để copy trọng số).
    Gradient có khớp đến từng chữ số không? Nếu không, vì sao? Gợi ý: kiểm
    tra xem bạn có đang chạy ở \mintinline{python}{torch.float64} không
    (\mintinline{python}{torch.set_default_dtype} để set mặc định).

  \item (Câu hỏi mở) Định lý universal approximation nói rằng một mạng
    truyền thẳng đủ rộng có thể xấp xỉ mọi hàm. Vậy vì sao ta cần RNN? Hãy
    trả lời bằng cách chỉ ra một bài toán mà mạng truyền thẳng \emph{không
    thể} giải được dù có bao nhiêu tham số đi nữa; và giải thích vì sao RNN
    giải được nó. Đây là câu hỏi chưa có trong kỳ thi nào; tôi thật sự muốn
    biết bạn nghĩ gì.
\end{bt}

\index{bài tập!đẩy xa hơn|)}%
